\documentclass[10pt,a4paper]{amsart}
\usepackage[margin=19mm]{geometry}
\usepackage{amsmath,amssymb,mathtools}
\usepackage[scr=rsfs,cal=euler]{mathalfa}
\usepackage{xfrac}
\usepackage[unicode,hidelinks,bookmarks=false]{hyperref}
\newcommand{\ie}{i.e.\ }
\newcommand{\cf}{cf.\ }
\newcommand{\resp}{respectively\ }
\newcommand{\Subsection}{\subsection}
\providecommand{\nobreakdash}{\nobreak\hbox{-}}
\setlength{\emergencystretch}{2em}
\multlinegap0pt
\usepackage{bm}
\usepackage{bbm}
\usepackage{dsfont}
\usepackage{xspace}
\usepackage{cancel}
\usepackage{mathtools}
\usepackage{amsthm}
\makeatletter
\@ifundefined{theorem}{\newtheorem{theorem}{Theorem}}{}
\@ifundefined{lemma}{\newtheorem{lemma}[theorem]{Lemma}}{}
\@ifundefined{proposition}{\newtheorem{proposition}[theorem]{Proposition}}{}
\makeatother
\makeatletter
\newcommand{\cG}{\mathcal{G}}
\expandafter\ifx\csname poc\endcsname\relax
\newenvironment{poc}{\begin{proof}[Proof of claim]\renewcommand*{\qedsymbol}{$\blacksquare$}}{\end{proof}}
\fi
\makeatother
\title[Geometric constructions for Ramsey-Tur\'an theory]{Geometric constructions for Ramsey-Tur\'an theory}
\author{Hong Liu, Christian Reiher, Maryam Sharifzadeh, Katherine Staden}
\date{Source: arXiv:2103.10423v2 (CC BY 4.0)}
\begin{document}
\maketitle

\noindent\emph{Source and licence.} Geometric constructions for Ramsey-Tur\'an theory. Version arXiv:2103.10423v2 (CC BY 4.0), distributed under the Creative Commons Attribution 4.0 International licence, as recorded in the arXiv licence metadata for this exact version and shown on the abstract page https://arxiv.org/abs/2103.10423v2. Full rights notice: licenses/arxiv-2103.10423-CC-BY-4.0.txt.\par\smallskip

\noindent\textit{Editorial scope.} Counted unit: the Lemma whose source label is \texttt{smallp} in the article (source0.tex lines 1824--1829, source label \texttt{smallp}), delivered together with the article's own complete original proof of it (source0.tex lines 1831--1952, one \texttt{proof} environment of 122 lines). No mathematics is written, repaired or re-derived: statement and proof are the article's own text line for line. Required context retained uncounted: prop-hero (lines 1756--1763). Every definition, display and label of those lines is retained unchanged. Omitted from this package and not redistributed: 1--1755, 1764--1823, 1830--1830, 1953--2200. Nothing inside a retained excerpt is omitted or altered: every retained body is delivered exactly as the article's lines are written, so no omission receipt of any kind accompanies this package. Citations: the retained excerpts carry no citation command at all, so no bibliography entry of the article is transcribed and no reference is redistributed. Reference adaptation: none was needed and none was made; every reference inside the delivered unit resolves inside it, so no unresolved reference or citation is produced. Printed slips kept literal: no printed word, symbol, hypothesis or number of the counted statement or of its proof was altered. Layout and class: the article's own class is \texttt{article} with packages including amssymb, amsmath, amsthm, fancyhdr, babel, enumitem, xcolor, hyperref, geometry, bbm, todonotes, graphicx, bm, amsfonts; the wrapper is the lane's pinned \texttt{amsart} editorial wrapper at 10pt on A4, which restates the theorem-like environments the retained text needs and adds the article's own macro definitions that the retained text needs (\texttt{\textbackslash cG}), with extra wrapper packages (bm, bbm, dsfont, xspace, cancel, mathtools). The delivered numbering is the unit's own. Assets: no retained span contains a figure environment, a graphics-inclusion command, a PDF-page inclusion, a table or a TikZ picture; no image file is copied into this package, the package declares no asset dependency and no image file is redistributed with it. Licence: version arXiv:2103.10423v2 of this article is distributed under the Creative Commons Attribution 4.0 International licence, as recorded in the arXiv licence metadata for this exact version and shown on the abstract page https://arxiv.org/abs/2103.10423v2. A full rights notice is delivered at licenses/arxiv-2103.10423-CC-BY-4.0.txt. Characters: every retained excerpt is pure ASCII, so the delivered engine, XeLaTeX with its default Latin Modern text font, reports no missing character.
\begin{proposition}\label{prop-hero}
	For every $p$-weighted graph $G=(V,w)$, there exists a set $K \subseteq V$ such that
	\begin{itemize}
	\item[(i)] $G[K] \in \cG_p(p|K|-\tilde{w}(K))$;
	\item[(ii)] we have $\gamma_K(y) \leq p-1$ for all $y \in K$ and $\gamma_K(x) \geq p$ for all $x \in V\setminus K$;
	\item[(iii)] if $x \in V\setminus K$ and $y \in K$, we have $\gamma_{K\setminus\{y\}}(x) \geq \gamma_K(y)$.
	\end{itemize}
\end{proposition}

\begin{lemma}\label{smallp}
	Let $p \in\lbrace 3,4\rbrace$ and let $t\in\mathbb{N}$.
	Let $G=(V,w)$ be a $p$-weighted $n$-vertex graph with
	$$\delta(G)>p \cdot \varrho_p^*(pt+2) \cdot n.$$
	Then there is $J \subseteq V$ such that $G[J] \in \cG_p(pt+2)$.
\end{lemma}

\begin{proof}
	Let $G=(V,w)$ be an $n$-vertex $p$-weighted graph with $\delta(G)>p\cdot\frac{(t-1)(2p-1)+1}{t(2p-1)+1}\cdot n$.	%We shall find a subgraph of $G$ which is in $\cG_p(pt+2)$.
	Let $K \subseteq V$ be the set obtained from Proposition~\ref{prop-hero}. We claim that
	\begin{equation}\label{eq-K}
		|K|\ge t+1.
	\end{equation}
	To see this, observe first that for each $y\in V$, by Proposition~\ref{prop-hero}(ii), 
	$$\sum_{x \in K}\tilde{w}(x,y)=\gamma_K(y)+p\cdot\mathbbm{1}_{\{y\in K\}}\ge p.$$
	Consequently,
	\begin{align}
		pn &\leq \sum_{x \in K,y \in V}\tilde{w}(x,y) =\sum_{x \in K}\Big(pn-\sum_{y\in V}w(x,y)\Big)\nonumber\\
		&=\sum_{x \in K}(pn-d_G(x)) \leq |K|(pn-\delta(G)) < \frac{(2p-1)pn|K|}{(2p-1)t+1} < \frac{pn|K|}{t},\label{eq:ricecake}
	\end{align}
	so $|K| \geq t+1$ as claimed.
	
	Now let $J \subseteq V$ be a set of vertices with enumeration 
	$$J = \lbrace x_1,\ldots,x_k,y_1,\ldots,y_r,z_1,\ldots,z_s\rbrace,$$
	equipped with a dominating extension $w$, such that $\lbrace x_1,\ldots,x_k\rbrace = K$, $w|_K$ has size $pk-\tilde{w}(K)$ and 
	\begin{itemize}
		\item  $w(y_i) \geq 2$ for all $i \in [r]$;
		\item  $w(z_i) \geq 1$ for all $i \in [s]$;
		\item $2r+s$ is maximal.
	\end{itemize}
	Such a pair $(J,w)$ does exist. Indeed, taking $r=s=0$, we see that $J=K$ has a dominating extension of size $pk-\tilde{w}(K)$ by Proposition~\ref{prop-hero}(i).
	
	By choice, $G[J] \in \mathcal{G}_p(pk-\tilde{w}(K)+2r+s)$.
	We shall argue that $G[J]$ is the desired subgraph in $\mathcal{G}_p(pt+2)$. Suppose otherwise, then
	$pk-\tilde{w}(K)+2r+s \leq pt+1$.
	Together with~\eqref{eq-K}, this implies
	\begin{equation}\label{double}
		(2p-1)k - 2\tilde{w}(K) + 4r+2s \leq (2p-1)t+1.
	\end{equation}
	In what follows, we write $\gamma := \gamma_K$. Define $\eta: J\rightarrow \mathbb{N}$ as follows:
	\begin{itemize}
		\item $\eta(x_i) := 2p-1-\gamma(x_i)$, for all $i \in [k]$;
		\item $\eta(y_i) := 4$, for all $i \in [r]$;
		\item $\eta(z_i) := 2$, for all $i \in [s]$.
	\end{itemize}
	Note that by Proposition~\ref{prop-hero}(ii) we have
	\begin{equation}\label{eq-etap}
	\eta(x) \geq p, \quad\text{for all }x \in K.
	\end{equation}
	Further define
	$$
	H(u) := \sum_{v \in J}\eta(v)\tilde{w}(u,v),\quad\text{for all }u \in V.
	$$
	Since $\sum_{i\in [k]}\gamma(x_i)=2\tilde{w}(K)$, we have as in~\eqref{eq:ricecake} that
	\begin{align*}
		\sum_{u \in V}H(u) &= \sum_{v \in J}\eta(v)\sum_{u \in V}\tilde{w}(u,v) \leq \Big(\sum_{v \in J}\eta(v)\Big)(pn-\delta(G))\\
		&< ((2p-1)k - 2\tilde{w}(K) + 4r+2s) \cdot \frac{p(2p-1)n}{(2p-1)t+1} \stackrel{(\ref{double})}{\leq} p(2p-1)n,
	\end{align*}
	implying that there exists a vertex $u_* \in V$ with $H(u_*) < p(2p-1)$.
	
	Suppose there is some $v \in K$ with $\tilde{w}(v,u_*) = p$.
	Then
	$$
	p(2p-1)>H(u_*)\ge \sum_{x' \in K\setminus\{v\}}\eta(x')\tilde{w}(x',u_*) + p\eta(v) \stackrel{(\ref{eq-etap})}{\ge} p\cdot\gamma_{K\setminus \{v\}}(u_*) + p(2p-1-\gamma(v)),
	$$
	implying that we in fact have $\gamma_{K\setminus\{v\}}(u_*) <  \gamma(v)$. Then Proposition~\ref{prop-hero}(iii) implies that $u_* \in K$.
	If $u_* \neq v$, then $\gamma(u_*) \geq \tilde{w}(v,u_*) = p$, contradicting Proposition~\ref{prop-hero}(ii).
	So $u_*=v$, and thus $\gamma_{K\setminus\{v\}}(u_*) = \gamma(v)$, and we have obtained a contradiction.
	Therefore,
	\begin{equation}\label{eq-K-linked}
		w(v,u_*) \geq 1,\quad\text{for all }v \in K.
	\end{equation}
	So $u_* \notin K$, since $w(u_*,u_*)=0$. Proposition~\ref{prop-hero}(ii) then implies that \begin{equation}\label{gammau*}
		\gamma(u_*) \geq p.
	\end{equation}
	Consequently, recalling that $\eta(x') \geq p$ for all $x' \in K$,
	\begin{equation}\label{eq-K-Hweight}
		\sum_{x' \in K}\eta(x')\tilde{w}(x',u_*) \geq p\gamma(u_*)\ge p^2.
	\end{equation}
	For every $y \in R := \lbrace y_1,\ldots,y_r\rbrace$, we have
	$$
	p(2p-1)>H(u_*)\ge 4\tilde{w}(y,u_*)+\sum_{x' \in K}\eta(x')\tilde{w}(x',u_*)\stackrel{\eqref{eq-K-Hweight}}{\ge} 4\tilde{w}(y,u_*)+p^2,
	$$
	whence $\tilde{w}(y,u_*)<\frac{p(p-1)}{4}<p$. So \begin{equation}\label{eq-R-linked}
		w(y,u_*) \geq 1,\quad\text{for all }y \in R.
	\end{equation}
	
	Suppose that $w(z,u_*) \geq 1$ for all $z \in S := \lbrace z_1,\ldots,z_s\rbrace$. Then $u_* \notin S$, by~\eqref{eq-K-linked} $u^* \notin K$ and by~\eqref{eq-R-linked}, $u_* \notin R$.
	So $u_* \notin J$ and moreover, setting $w(u_*)=1$ gives a dominating extension to include $J \cup \{u_*\}$.
	Thus we can add $u_*$ to $S$, contradicting the maximality of $2r+s$.
	So $S' := \{z \in S: \tilde{w}(z,u_*)=p\} \neq \varnothing$.
	Then  we have
	$$p(2p-1)>H(u_*)\ge \sum_{z \in S'}\eta(z)\tilde{w}(z,u_*)+\sum_{x' \in K}\eta(x')\tilde{w}(x',u_*)\stackrel{\eqref{eq-K-Hweight}}{\ge} 2|S'|p+p\gamma(u_*), $$
	implying together with~(\ref{gammau*}) that
	\begin{equation}\label{eq:S'}
		|S'|<\frac{2p-1-\gamma(u_*)}{2} \le \frac{p-1}{2}.
	\end{equation}
	If $p=3$, then $S'$ is empty, a contradiction.
	So from now on assume $p=4$.
	Then $|S'|=1$, i.e.~there is a unique $z_* \in S$ with $\tilde{w}(z_*,u_*)=4$, and 
	\begin{equation}\label{Sweight}
		w(z,u_*)\geq 1,\quad\text{for all }z \in S\setminus\lbrace z_*\rbrace.
	\end{equation}
	The first inequality in~(\ref{eq:S'}) implies that $\gamma(u_*) \leq 4$.
	Now~(\ref{gammau*}) implies that in fact $$\gamma(u_*)=4.$$
	Also, from
	$$
	28 > H(u_*) \geq 2\tilde{w}(z_*,u_*)+ 4\sum_{y \in R}\tilde{w}(y,u_*) +\sum_{x' \in K}\eta(x')\tilde{w}(x',u_*)\stackrel{\eqref{eq-K-Hweight}}{\ge} 8+4\sum_{y \in R}\tilde{w}(y,u_*)+16,
	$$
	we deduce that $\sum_{y \in R}\tilde{w}(y,u_*)=0$. In other words,
	\begin{equation}\label{Rweight}
		w(y,u_*) =4, \quad\text{for all }y \in R.
	\end{equation}
	
	Let $T := \lbrace x \in K: \tilde{w}(x,u_*) > 0\rbrace$. By definition, $ \sum_{x \in T}\tilde{w}(x,u_*)=\gamma(u_*) =4$, implying, together with~\eqref{eq-K-linked}, that $2 \leq |T| \leq 4$.
	If $|T| \geq 3$, then the multiset $\lbrace w(x,u_*) : x \in T\rbrace$ recording the weights from $u_*$ to $T$ is either $ \lbrace 3,3,3,3\rbrace$, or  $\lbrace 2,3,3\rbrace$. In particular, by~(\ref{Rweight}), $w(y,u_*) \geq 2$ for all $y \in K \cup R$ with equality at most once.
	Together with~(\ref{Sweight}), this implies that
	we can delete $z_*$ from $S$ and add $u_*$ to $R$ to obtain a set, $J\cup\{u_*\}\setminus\{z_*\}$, having a dominating extension with larger $2r+s$, a contradiction.
	
	Thus we may assume that $|T| = 2$.
	Let $T =: \lbrace a,b\rbrace$ and $\alpha := \tilde{w}(a,u_*)$ and $\beta := \tilde{w}(b,u_*)$.
	Then $\alpha+\beta=\gamma(u_*)=4$; and Proposition~\ref{prop-hero}(iii) implies that 
	$\beta = \gamma_{T\setminus\{a\}}(u_*) = \gamma_{K\setminus\{a\}}(u_*) \geq \gamma(a)$
	and similarly $\alpha \geq \gamma(b)$. We then arrive at the final contradiction:
	\begin{align*}
		28>H(u_*) &\geq \sum_{x' \in \{a,b,z_*\}}\eta(x')\tilde{w}(x',u_*) \ge \alpha(7-\gamma(a)) + \beta(7-\gamma(b)) + 2\cdot 4\\ &\ge \alpha(7-\beta) + \beta(7-\alpha) + 8\ge  8 + 7(\alpha+\beta)-\frac{1}{2}(\alpha+\beta)^2=28,
	\end{align*}
	completing the proof.
\end{proof}

\end{document}
